\documentclass[12pt,a4paper]{article}

% ---- Encoding & Font ----
\usepackage{iftex}
\ifPDFTeX
  \usepackage[utf8]{inputenc}
  \usepackage[T1]{fontenc}
\else
  \usepackage{fontspec}  % XeTeX/LuaTeX (tectonic): Unicode nativo (·, —, ¿, ª, §)
\fi
\usepackage[spanish]{babel}
\usepackage{csquotes}

% ---- Page layout ----
\usepackage[top=2.5cm, bottom=2.5cm, left=2.5cm, right=2.5cm]{geometry}
\usepackage{setspace}
\onehalfspacing

% ---- Math ----
\usepackage{amsmath, amssymb, amsthm}
\usepackage{mathtools}
\usepackage{physics}
\usepackage{esint}  % \oiint

% ---- Graphics ----
\usepackage{graphicx}
\usepackage{tikz}
\usepackage{float}
\usepackage{caption}

% ---- Tables ----
\usepackage{array}
\usepackage{booktabs}
\usepackage{tabularx}
\usepackage{colortbl}

% ---- Colours ----
\usepackage{xcolor}
\definecolor{notebg}{HTML}{FFF3CD}
\definecolor{noteborder}{HTML}{FFC107}
\definecolor{boxred}{HTML}{F8D7DA}
\definecolor{boxredborder}{HTML}{F5C6CB}

% ---- Boxes ----
\usepackage{mdframed}
\usepackage{tcolorbox}
\tcbuselibrary{most}

\newtcolorbox{keybox}[1][]{
  colback=blue!5,
  colframe=blue!70!black,
  arc=4pt,
  boxrule=1pt,
  fonttitle=\bfseries,
  title={#1}
}

\newtcolorbox{warnbox}[1][]{
  colback=boxred,
  colframe=boxredborder,
  coltitle=black!75,
  arc=4pt,
  boxrule=1pt,
  fonttitle=\bfseries,
  title={#1}
}

\newtcolorbox{notebox}[1][]{
  colback=notebg,
  colframe=noteborder,
  coltitle=black!75,
  arc=4pt,
  boxrule=1pt,
  fonttitle=\bfseries,
  title={#1}
}

% ---- Hyperlinks & TOC ----
\usepackage[hidelinks]{hyperref}
\usepackage{bookmark}

% ---- Enumerate ----
\usepackage{enumitem}

% ---- Miscellaneous ----
\usepackage{icomma}
\usepackage{siunitx}
\usepackage{textcomp}
\usepackage[bottom]{footmisc}

% ---- Diagramas (gráficas y esquemas) ----
\usepackage{pgfplots}
\pgfplotsset{compat=1.18}
\usetikzlibrary{babel}  % babel español activa `>`; esto evita romper las flechas `->`
% courses/ElecMag2024/Clases/assets/tikzstyles.tex
% ---------------------------------------------------------------------------
% Estilos compartidos para los diagramas del curso Electromagnetismo (2024).
% Fuente única del "look" visual (colores, grosores, escalas) para que todas
% las figuras (TikZ / pgfplots / CircuiTikz) sean consistentes entre clases.
%
% Es una COPIA de courses/Fisica3-2015/Clases/assets/tikzstyles.tex: mismo
% dominio, misma paleta. Copia y no symlink para que cada curso sea
% autocontenido y la paleta de Física III pueda evolucionar aparte.
%
% Uso: la clase debe cargar antes `tikz` y `pgfplots` (y `circuitikz` si la
%      clase tiene circuitos), y luego % courses/ElecMag2024/Clases/assets/tikzstyles.tex
% ---------------------------------------------------------------------------
% Estilos compartidos para los diagramas del curso Electromagnetismo (2024).
% Fuente única del "look" visual (colores, grosores, escalas) para que todas
% las figuras (TikZ / pgfplots / CircuiTikz) sean consistentes entre clases.
%
% Es una COPIA de courses/Fisica3-2015/Clases/assets/tikzstyles.tex: mismo
% dominio, misma paleta. Copia y no symlink para que cada curso sea
% autocontenido y la paleta de Física III pueda evolucionar aparte.
%
% Uso: la clase debe cargar antes `tikz` y `pgfplots` (y `circuitikz` si la
%      clase tiene circuitos), y luego % courses/ElecMag2024/Clases/assets/tikzstyles.tex
% ---------------------------------------------------------------------------
% Estilos compartidos para los diagramas del curso Electromagnetismo (2024).
% Fuente única del "look" visual (colores, grosores, escalas) para que todas
% las figuras (TikZ / pgfplots / CircuiTikz) sean consistentes entre clases.
%
% Es una COPIA de courses/Fisica3-2015/Clases/assets/tikzstyles.tex: mismo
% dominio, misma paleta. Copia y no symlink para que cada curso sea
% autocontenido y la paleta de Física III pueda evolucionar aparte.
%
% Uso: la clase debe cargar antes `tikz` y `pgfplots` (y `circuitikz` si la
%      clase tiene circuitos), y luego \input{../assets/tikzstyles.tex}
% (compilar desde el directorio de la clase para que la ruta relativa resuelva).
% ---------------------------------------------------------------------------

% ---- Paleta (alineada con las cajas del preámbulo: keybox/notebox/warnbox) ----
\definecolor{figblue}{HTML}{1F4E79}   % azul principal (líneas, keybox)
\definecolor{figred}{HTML}{C0392B}    % rojo de énfasis (warnbox)
\definecolor{figamber}{HTML}{B8860B}  % ámbar (notebox)
\definecolor{figgray}{HTML}{555555}   % auxiliares, ejes, guías

% ---- CircuiTikz: escala y grosor de los componentes ----
% Guarda \ifdefined: la electrostática (Clases 1-14) no carga `circuitikz`, y
% sin la guarda el \input revienta con `Undefined control sequence \ctikzset`.
% Cuando el curso llegue a circuitos, el bloque se activa solo.
\ifdefined\ctikzset
  \ctikzset{
    bipoles/thickness=1,
    resistors/scale=0.9,
    capacitors/scale=0.9,
  }
\fi

% ---- TikZ: estilos reutilizables ----
% Nota: las puntas se declaran como `-{Latex[...]}` y no como `->`. La forma
% con llaves NO contiene un `>` literal, así que es inmune al `>` activo que
% introduce babel español (ver CLAUDE.md §7.3.3).
\usetikzlibrary{arrows.meta,calc,angles,quotes}

\tikzset{
  figline/.style={figblue, line width=0.9pt},
  figaux/.style={figgray, dashed, line width=0.6pt},
  % vectores
  figvec/.style={figblue, line width=0.9pt, -{Latex[length=2.2mm,width=1.6mm]}},
  figvecres/.style={figred, line width=1.3pt, -{Latex[length=2.6mm,width=1.9mm]}},
  figvecaux/.style={figgray, line width=0.7pt, -{Latex[length=1.8mm,width=1.3mm]}},
  figaxis/.style={figgray, line width=0.6pt, -{Latex[length=2mm,width=1.4mm]}},
  % cargas puntuales
  figcarga/.style={circle, inner sep=0pt, minimum size=5.4pt, draw=none},
  figpos/.style={figcarga, fill=figred},
  figneg/.style={figcarga, fill=figblue},
  figneutra/.style={figcarga, fill=figgray},
  % etiquetas
  figlbl/.style={font=\small},
  figlblaux/.style={font=\footnotesize, figgray},
}

% ---- pgfplots: estilo base de las gráficas ----
\pgfplotsset{
  emfig/.style={
    width=11cm, height=6.5cm,
    axis lines=left,
    xlabel style={font=\small},
    ylabel style={font=\small},
    tick label style={font=\footnotesize},
    every axis plot/.append style={line width=1pt},
    grid=major,
    grid style={figgray!20},
    legend cell align=left,
    legend style={font=\footnotesize, draw=figgray!40},
  },
}

% (compilar desde el directorio de la clase para que la ruta relativa resuelva).
% ---------------------------------------------------------------------------

% ---- Paleta (alineada con las cajas del preámbulo: keybox/notebox/warnbox) ----
\definecolor{figblue}{HTML}{1F4E79}   % azul principal (líneas, keybox)
\definecolor{figred}{HTML}{C0392B}    % rojo de énfasis (warnbox)
\definecolor{figamber}{HTML}{B8860B}  % ámbar (notebox)
\definecolor{figgray}{HTML}{555555}   % auxiliares, ejes, guías

% ---- CircuiTikz: escala y grosor de los componentes ----
% Guarda \ifdefined: la electrostática (Clases 1-14) no carga `circuitikz`, y
% sin la guarda el \input revienta con `Undefined control sequence \ctikzset`.
% Cuando el curso llegue a circuitos, el bloque se activa solo.
\ifdefined\ctikzset
  \ctikzset{
    bipoles/thickness=1,
    resistors/scale=0.9,
    capacitors/scale=0.9,
  }
\fi

% ---- TikZ: estilos reutilizables ----
% Nota: las puntas se declaran como `-{Latex[...]}` y no como `->`. La forma
% con llaves NO contiene un `>` literal, así que es inmune al `>` activo que
% introduce babel español (ver CLAUDE.md §7.3.3).
\usetikzlibrary{arrows.meta,calc,angles,quotes}

\tikzset{
  figline/.style={figblue, line width=0.9pt},
  figaux/.style={figgray, dashed, line width=0.6pt},
  % vectores
  figvec/.style={figblue, line width=0.9pt, -{Latex[length=2.2mm,width=1.6mm]}},
  figvecres/.style={figred, line width=1.3pt, -{Latex[length=2.6mm,width=1.9mm]}},
  figvecaux/.style={figgray, line width=0.7pt, -{Latex[length=1.8mm,width=1.3mm]}},
  figaxis/.style={figgray, line width=0.6pt, -{Latex[length=2mm,width=1.4mm]}},
  % cargas puntuales
  figcarga/.style={circle, inner sep=0pt, minimum size=5.4pt, draw=none},
  figpos/.style={figcarga, fill=figred},
  figneg/.style={figcarga, fill=figblue},
  figneutra/.style={figcarga, fill=figgray},
  % etiquetas
  figlbl/.style={font=\small},
  figlblaux/.style={font=\footnotesize, figgray},
}

% ---- pgfplots: estilo base de las gráficas ----
\pgfplotsset{
  emfig/.style={
    width=11cm, height=6.5cm,
    axis lines=left,
    xlabel style={font=\small},
    ylabel style={font=\small},
    tick label style={font=\footnotesize},
    every axis plot/.append style={line width=1pt},
    grid=major,
    grid style={figgray!20},
    legend cell align=left,
    legend style={font=\footnotesize, draw=figgray!40},
  },
}

% (compilar desde el directorio de la clase para que la ruta relativa resuelva).
% ---------------------------------------------------------------------------

% ---- Paleta (alineada con las cajas del preámbulo: keybox/notebox/warnbox) ----
\definecolor{figblue}{HTML}{1F4E79}   % azul principal (líneas, keybox)
\definecolor{figred}{HTML}{C0392B}    % rojo de énfasis (warnbox)
\definecolor{figamber}{HTML}{B8860B}  % ámbar (notebox)
\definecolor{figgray}{HTML}{555555}   % auxiliares, ejes, guías

% ---- CircuiTikz: escala y grosor de los componentes ----
% Guarda \ifdefined: la electrostática (Clases 1-14) no carga `circuitikz`, y
% sin la guarda el \input revienta con `Undefined control sequence \ctikzset`.
% Cuando el curso llegue a circuitos, el bloque se activa solo.
\ifdefined\ctikzset
  \ctikzset{
    bipoles/thickness=1,
    resistors/scale=0.9,
    capacitors/scale=0.9,
  }
\fi

% ---- TikZ: estilos reutilizables ----
% Nota: las puntas se declaran como `-{Latex[...]}` y no como `->`. La forma
% con llaves NO contiene un `>` literal, así que es inmune al `>` activo que
% introduce babel español (ver CLAUDE.md §7.3.3).
\usetikzlibrary{arrows.meta,calc,angles,quotes}

\tikzset{
  figline/.style={figblue, line width=0.9pt},
  figaux/.style={figgray, dashed, line width=0.6pt},
  % vectores
  figvec/.style={figblue, line width=0.9pt, -{Latex[length=2.2mm,width=1.6mm]}},
  figvecres/.style={figred, line width=1.3pt, -{Latex[length=2.6mm,width=1.9mm]}},
  figvecaux/.style={figgray, line width=0.7pt, -{Latex[length=1.8mm,width=1.3mm]}},
  figaxis/.style={figgray, line width=0.6pt, -{Latex[length=2mm,width=1.4mm]}},
  % cargas puntuales
  figcarga/.style={circle, inner sep=0pt, minimum size=5.4pt, draw=none},
  figpos/.style={figcarga, fill=figred},
  figneg/.style={figcarga, fill=figblue},
  figneutra/.style={figcarga, fill=figgray},
  % etiquetas
  figlbl/.style={font=\small},
  figlblaux/.style={font=\footnotesize, figgray},
}

% ---- pgfplots: estilo base de las gráficas ----
\pgfplotsset{
  emfig/.style={
    width=11cm, height=6.5cm,
    axis lines=left,
    xlabel style={font=\small},
    ylabel style={font=\small},
    tick label style={font=\footnotesize},
    every axis plot/.append style={line width=1pt},
    grid=major,
    grid style={figgray!20},
    legend cell align=left,
    legend style={font=\footnotesize, draw=figgray!40},
  },
}


% ---- Title ----
\title{\textbf{Clase 17: Ley de Biot y Savart} \\[0.3em]
        \large{Ausencia de monopolos magnéticos, símbolo de Levi-Civita y ejemplos de aplicación}}
\author{Transcripto y expandido --- Curso Electromagnetismo (OpenFING)}
\date{}

\begin{document}

\maketitle
\thispagestyle{empty}
\newpage

\tableofcontents
\newpage

\section{De las fuerzas a las fuentes del campo}

La clase anterior trató las \textbf{fuerzas} que cargas en movimiento y corrientes
sienten en un campo magnético. Esta clase y la siguiente tratan lo otro: cómo
corrientes y cargas en movimiento \textbf{generan} campos magnéticos. La primera ley
es la de \textbf{Biot y Savart}, conocida de Física III, a la que ahora se le va a
sacar más jugo.

Contexto histórico, tal como lo cuenta el docente:

\begin{itemize}
  \item El experimento mencionado la clase pasada (en el VTT, «Hertz»; ver la nota de
  la Clase 16) era \textbf{cualitativo}: mostraba que una corriente produce efectos
  magnéticos, sin números.
  \item El primero en estudiar \textbf{experimental y cuantitativamente} cómo las
  corrientes generan efectos magnéticos fue \textbf{Ampère}, alrededor de 1820.
  \item La expresión matemática general la dieron luego Biot y Savart.
\end{itemize}

\begin{notebox}
El docente dice que Biot y Savart no trabajaron juntos, que uno era inglés y
el otro francés, y que por eso los manuales franceses y británicos nombran la
ley distinto. Se registra como dicho en clase; no coincide con la historia
habitual (ambos eran franceses y publicaron juntos), así que conviene no
tomarlo como dato.
\end{notebox}

\section{Ley de Biot y Savart}

\subsection{Campo de una carga en movimiento}

Punto de partida, de la clase pasada: una carga $q_1$ en la posición $\vec r_1$
con velocidad $\vec v_1$ produce en $\vec r_2$ el campo

\[
\vec B(\vec r_2) = \frac{\mu_0}{4\pi}\,q_1\,
\frac{\vec v_1\times(\vec r_2-\vec r_1)}{|\vec r_2-\vec r_1|^3} .
\]

Pero el efecto de cargas individuales a velocidades normales es prácticamente
indetectable; lo que interesa son las corrientes en sistemas neutros.

\subsection{Densidad de corriente distribuida}

Se toma un conjunto de portadores con $N$ cargas por unidad de volumen. En un
volumen pequeño $dV$ en torno de $\vec r_1$ cada carga tiene su propia
velocidad, pero se las reemplaza a todas por la \textbf{velocidad de deriva}. El
campo que ese volumencito produce en $\vec r_2$ es el de una carga por el
número de cargas, $N\,dV$:

\[
d\vec B(\vec r_2) = \frac{\mu_0}{4\pi}\;N\,dV\;q_1\,
\frac{\vec v_1\times(\vec r_2-\vec r_1)}{|\vec r_2-\vec r_1|^3} .
\]

(Se supone un solo tipo de portador, como en un metal; si no, se suma sobre los
tipos.) Usando $\vec J(\vec r_1) = N\,q_1\,\vec v_1$ e integrando sobre todo el
volumen donde hay corriente:

\begin{keybox}
\[
\boxed{\ \vec B(\vec r_2) = \frac{\mu_0}{4\pi}\int_V
\frac{\vec J(\vec r_1)\times(\vec r_2-\vec r_1)}{|\vec r_2-\vec r_1|^3}\,dV_1\ }
\]
\end{keybox}

\textbf{Ley de Biot y Savart para una densidad de corriente.}

\subsection{Corriente en un cable}

Lo habitual es que la corriente esté concentrada en cables. El pasaje se hizo
con las cuentas precisas la clase pasada (sección del cable, etc.); ahora se
usa como \textbf{regla mnemotécnica}, dimensionalmente correcta:

\[
\vec J\,dV \;\longrightarrow\; I\,d\vec l .
\]

La integral de volumen pasa a ser una integral de línea sobre la curva $C$ del
circuito (normalmente cerrada), y la corriente, que es la misma en todo el
cable, sale afuera:

\begin{keybox}
\[
\boxed{\ \vec B(\vec r_2) = \frac{\mu_0 I}{4\pi}\oint_C
\frac{d\vec l_1\times(\vec r_2-\vec r_1)}{|\vec r_2-\vec r_1|^3}\ }
\]
\end{keybox}

\textbf{Ley de Biot y Savart para corrientes en cables}, la forma de Física III y la
que históricamente se midió: es muy difícil medir una distribución, lo que se
mide es corriente en cables.

\begin{warnbox}[Dominio de validez]
Para corrientes que no varían mucho en el tiempo, es
la expresión general del campo magnético producido por una corriente y
contiene, en principio, toda la información. Cuando las corrientes varían
mucho aparecen efectos nuevos, que se verán más adelante.
\end{warnbox}

\section{No hay monopolos magnéticos}

\subsection{Enunciado}

Las líneas de campo magnético se cierran sobre sí mismas o van al infinito, pero
\textbf{no terminan en un punto}. Es distinto del campo eléctrico, cuyas líneas salen
de las cargas positivas y terminan en las negativas. La ley de Gauss dice
$\div \vec E = \rho/\varepsilon_0$ en el vacío; si no hay monopolos
magnéticos, el equivalente de la densidad de carga es cero:

\begin{keybox}
\[
\boxed{\ \div \vec B = 0\ }
\]
\end{keybox}

En Física III se escribía en forma integral: el flujo de $\vec B$ a través de
cualquier superficie cerrada es nulo. Por el teorema de la divergencia son
equivalentes, y la forma diferencial termina siendo más cómoda.

\begin{notebox}[Es un hecho experimental]
Aunque acá se deduce de Biot y Savart, tiene que
verificarse en el laboratorio. Los físicos se preguntan desde el siglo XIX si
no habrá monopolos magnéticos: los teóricos «desearíamos» que los hubiera
---explicarían varias cosas---, y los experimentales se empecinan en que no hay.
Muchas teorías que generalizan el electromagnetismo predicen cierta densidad
de monopolos en el universo, y fueron descartadas por eso.
\end{notebox}

\subsection{Demostración}

Se trabaja con la forma de volumen, que es menos singular (cualquier cable,
mirado de cerca, tiene una distribución). Se toma la divergencia respecto de
las coordenadas del punto de campo, $\vnabla_2$ (derivadas en $x_2, y_2, z_2$).

Hace falta la divergencia de un producto vectorial:

\[
\div (\vec F\times\vec G) = \vec G\cdot(\curl \vec F) - \vec F\cdot(\curl \vec G) .
\]

\begin{notebox}
En estas fórmulas todo es «medio obvio» ---es como la derivada de un producto y,
para que cierren los vectores, uno de los factores tiene que ser un
rotacional y el otro entrar con producto escalar--- \textbf{salvo el signo}, porque
el producto vectorial no es conmutativo. El truco de §4 es para no
equivocarse en eso.
\end{notebox}

Con $\vec F = \vec J(\vec r_1)$ y
$\vec G = (\vec r_2-\vec r_1)/|\vec r_2-\vec r_1|^3$:

\begin{enumerate}
  \item $\vec J(\vec r_1)$ no depende de $\vec r_2$, así que
  $\vnabla_2\times\vec J(\vec r_1) = 0$: el primer término se anula.
  \item Queda
  \[
  \vnabla_2\cdot\vec B(\vec r_2) = -\frac{\mu_0}{4\pi}\int_V \vec J(\vec r_1)\cdot
  \left[\vnabla_2\times\frac{\vec r_2-\vec r_1}{|\vec r_2-\vec r_1|^3}\right]dV_1 .
  \]
  \item El corchete es el rotacional de algo proporcional al \textbf{campo eléctrico de
  una carga puntual}, que es cero; esa cuenta ya se hizo al estudiar Coulomb.
\end{enumerate}

Por lo tanto $\vnabla_2\cdot\vec B(\vec r_2) = 0$. $\blacksquare$

Son dos las leyes fundamentales del campo magnético: la ausencia de monopolos,
que sale de tomar la \textbf{divergencia}, y la ley de Ampère, que saldrá la próxima
clase de tomar el \textbf{rotacional}.

\section{Paréntesis matemático: el símbolo de Levi-Civita}

Al docente le molesta ir introduciendo fórmulas de cálculo vectorial «sacadas
de la manga». La comparación es con la trigonometría, que en Física III reduce a
una única fórmula (la de Euler o de De Moivre). Acá hace lo mismo: \textbf{todas} las
fórmulas diferenciales del cálculo vectorial salen de una herramienta.

\begin{notebox}
El docente lo atribuye a «dos matemáticos italianos, Levi y Civita». Es una
sola persona: Tullio Levi-Civita.
\end{notebox}

\subsection{Definición}

$\varepsilon_{ijk}$ es un símbolo con tres índices, cada uno con valores 1, 2
o 3 (las coordenadas cartesianas $x$, $y$, $z$), definido por dos propiedades:

\[
\varepsilon_{123} = 1,
\qquad
\varepsilon_{ijk} = -\varepsilon_{jik} = -\varepsilon_{kji}
\quad\text{(completamente antisimétrico).}
\]

Al intercambiar dos índices cualesquiera cambia de signo. Iterando, todos los
casos se reducen a esos.

\begin{center}
\renewcommand{\arraystretch}{1.35}
\small
\begin{tabularx}{\textwidth}{>{\raggedright\arraybackslash}X>{\raggedright\arraybackslash}X>{\raggedright\arraybackslash}X}
\toprule
\textbf{Caso} & \textbf{Valor} & \textbf{Por qué} \\
\midrule
Índice repetido, p. ej. $\varepsilon_{113}$ & $0$ & Intercambiando los dos iguales es igual a menos sí mismo \\
Permutación par, p. ej. $\varepsilon_{231}$ & $+1$ & Dos intercambios desde 123 \\
Permutación impar, p. ej. $\varepsilon_{321}$, $\varepsilon_{132}$ & $-1$ & Un intercambio desde 123 \\
\bottomrule
\end{tabularx}
\end{center}

\subsection{Producto vectorial y rotacional}

Sirve porque es una manera muy cómoda de escribir productos vectoriales. Todas
las sumas van de 1 a 3:

\begin{keybox}
\[
\boxed{\ (\vec a\times\vec b)_i = \sum_{j,k}\varepsilon_{ijk}\,a_j\,b_k\ }
\]
\end{keybox}

El índice $i$ no está sumado: es un dato.

\textbf{Verificación para $i=3$.} $j$ y $k$ no pueden valer 3 ni ser iguales entre
sí (el símbolo se anularía), así que de toda la suma sobreviven dos términos:

\[
(\vec a\times\vec b)_3 = \varepsilon_{312}\,a_1 b_2 + \varepsilon_{321}\,a_2 b_1
= a_1 b_2 - a_2 b_1,
\]
que es la componente $z$ conocida. La fórmula no tiene ambigüedad: fijado $i$,
no hay opción para los otros dos.

Del mismo modo, para el rotacional:

\[
(\curl \vec v)_i = \sum_{j,k}\varepsilon_{ijk}\,\partial_j v_k .
\]

La ventaja: $\partial_j$ es una \textbf{derivada común y silvestre} respecto de
$x_j$, no un operador vectorial; se la manipula con las reglas usuales.

\begin{notebox}
Está hecho en cartesianas, pero una fórmula \textbf{vectorial} probada en un
sistema de coordenadas vale en cualquiera. Hay generalizaciones del símbolo a
otros sistemas, más complicadas.
\end{notebox}

\subsection{Ejemplo: divergencia de un producto vectorial}

Se prueba la fórmula usada en §3.2. La divergencia es derivar cada componente
respecto de su coordenada y sumar:

\[
\div (\vec F\times\vec G) = \sum_i \partial_i(\vec F\times\vec G)_i
= \sum_{i,j,k}\varepsilon_{ijk}\,\partial_i\!\left(F_j\,G_k\right),
\]
donde $\varepsilon_{ijk}$, que es constante, salió de la derivada. Por la regla
del producto:

\[
= \sum_{i,j,k}\varepsilon_{ijk}\left[(\partial_i F_j)\,G_k + F_j\,(\partial_i G_k)\right].
\]

\textbf{Primer término.} Sumando en $i$ y $j$, los índices están en el orden
correcto y, por permutación circular,
$\sum_{i,j}\varepsilon_{ijk}\,\partial_i F_j$ es la componente $k$ del
rotacional de $\vec F$. Queda $\sum_k (\curl \vec F)_k\,G_k$, que es la
definición misma de producto escalar:

\[
\vec G\cdot(\curl \vec F).
\]

\textbf{Segundo término.} Acá los índices \textbf{no} están en el orden correcto: la
derivada y la componente de $\vec G$ quedan con el primero y el tercero, con el
de $\vec F$ en el medio. Se usa la antisimetría para intercambiar $i$ y $j$, lo
que agrega un signo:

\[
\sum_{i,j,k}\varepsilon_{ijk}\,F_j\,\partial_i G_k
= -\sum_{i,j,k}\varepsilon_{jik}\,F_j\,\partial_i G_k
= -\,\vec F\cdot(\curl \vec G).
\]

Sumando:

\begin{keybox}
\[
\boxed{\ \div (\vec F\times\vec G)
= \vec G\cdot(\curl \vec F) - \vec F\cdot(\curl \vec G)\ }
\]
\end{keybox}

\begin{notebox}[Recomendación del docente]
«Esto es como ver al profesor de gimnasia hacer
gimnasia: no sirve para nada.» Hay que rehacer la cuenta paso a paso. No hay
nada conceptualmente difícil; es operatoria que se automatiza. Quien no
quiera usarlo puede seguir calculando «a pedal», componente a componente.
\end{notebox}

\section{Ejemplos de aplicación de Biot y Savart}

\subsection{Alcance: corrientes dadas}

Tramo de repaso de Física III, hecho ahora «por fuerza bruta» con productos
vectoriales explícitos ---no porque haga falta, sino como ejercitación, para que
el método se pueda generalizar.

Si las corrientes son \textbf{conocidas}, el problema magnético se reduce a hacer una
integral (analítica o, la mayoría de las veces, numérica).

\begin{notebox}
Eso no es toda la historia. Igual que en electrostática una esfera se
polarizaba y había que resolver a la vez la distribución de carga, las
propias corrientes pueden reaccionar al campo magnético. Esos problemas se
tratan con métodos parecidos a los de Laplace, más adelante.
\end{notebox}

\subsection{Cable recto muy largo}

Cable sobre el eje $x$ con corriente $I$ constante; punto de observación sobre
el eje $y$ a distancia $a$ del origen.

\[
\vec r_1 = x\,\hat\imath,\qquad \vec r_2 = a\,\hat\jmath,\qquad
d\vec l_1 = dx\,\hat\imath .
\]

\begin{figure}[H]
  \centering
  % Clase 17 — cable recto muy largo sobre el eje x; punto de observación sobre
% el eje y a distancia a. theta es el ángulo, en el elemento de corriente,
% entre el cable y la recta al punto: sin(theta) = a / sqrt(a^2 + x^2).
% Dibujado con x = -a (theta = 45 grados). El campo sale del papel (k).
\begin{tikzpicture}[scale=1]

  % ejes
  \draw[figaxis] (0,-0.5) -- (0,3.3);
  \node[figlbl] at (0,3.55) {$y$};
  \node[figlbl] at (5.45,0) {$x$};

  % cable
  \draw[figblue, line width=1.6pt, -{Latex[length=2.6mm,width=1.9mm]}] (-5,0) -- (5.1,0);
  \draw[figvecres] (2.6,0.3) -- (3.5,0.3);
  \node[figlbl, text=figred] at (3.05,0.62) {$I$};

  % elemento de corriente en x = -a
  \draw[figred, line width=2.4pt] (-2.6,0) -- (-2.0,0);
  \node[figlbl, text=figred] at (-2.3,-0.35) {$dx$};

  % vector al punto
  \draw[figaux, -{Latex[length=2mm,width=1.4mm]}] (-2.3,0) -- (-0.06,2.24);
  \node[figlbl] at (-1.85,1.3) {$\vec r_2-\vec r_1$};

  % angulo
  \draw[figgray, line width=0.6pt] (-1.5,0) arc (0:45:0.8);
  \node[figlbl] at (-1.15,0.33) {$\theta$};

  % punto de observacion y campo saliente
  \fill[figblue] (0,2.3) circle (1.6pt);
  \node[figlbl, anchor=west] at (0.2,1.25) {$a$};
  \draw[figred, line width=0.9pt] (0.55,2.3) circle (0.17);
  \fill[figred] (0.55,2.3) circle (0.05);
  \node[figlbl, text=figred, anchor=west] at (0.8,2.3) {$\vec B$};

  \node[figlblaux] at (-0.22,-0.3) {$O$};

\end{tikzpicture}

  \caption{Cable sobre el eje $x$ y punto de observación sobre el eje $y$. El ángulo $\theta$ va de $0$ a $\pi$ cuando el elemento de corriente recorre el cable de $-\infty$ a $+\infty$. El campo sale del papel.}
\end{figure}

El producto vectorial es inmediato ($\hat\imath\times\hat\imath = 0$,
$\hat\imath\times\hat\jmath = \hat k$):

\[
d\vec l_1\times(\vec r_2-\vec r_1) = dx\,\hat\imath\times(a\,\hat\jmath - x\,\hat\imath)
= a\,dx\,\hat k .
\]

\[
\vec B(\vec r_2) = \frac{\mu_0 I}{4\pi}\,\hat k\int_{-\infty}^{+\infty}
\frac{a\,dx}{\left(a^2+x^2\right)^{3/2}} .
\]

El problema físico está resuelto; falta una integral. En Física III se vieron
otros trucos (seno hiperbólico, partes); acá se usa un \textbf{cambio a variable
angular}. Sea $\theta$ el ángulo, en el elemento de corriente, entre el cable
y la recta al punto de observación:

\[
\sin\theta = \frac{a}{\sqrt{a^2+x^2}},
\qquad
x = -\,a\,\frac{\cos\theta}{\sin\theta}
\quad\Longrightarrow\quad
dx = \frac{a}{\sin^2\theta}\,d\theta .
\]

Cuando $x\to-\infty$, $\theta\to0$; cuando $x\to+\infty$, $\theta\to\pi$ (el
docente discute los límites un rato con los estudiantes hasta fijarlos). Como
$\left(a^2+x^2\right)^{-3/2} = \sin^3\theta/a^3$:

\[
\vec B = \frac{\mu_0 I}{4\pi a}\,\hat k\int_0^{\pi}\sin\theta\,d\theta
= \frac{\mu_0 I}{4\pi a}\,\hat k\,\Big[-\cos\theta\Big]_0^{\pi}
\]

\begin{keybox}
\[
\boxed{\ \vec B = \frac{\mu_0 I}{2\pi a}\,\hat k\ }
\]
\end{keybox}

Es el resultado conocido, que por Ampère sale en dos líneas. Se calculó en un
punto particular; al rotar el punto en torno del cable el campo rota con él:
las líneas son curvas cerradas alrededor del cable, en planos transversales.

\subsection{Espira circular, campo sobre el eje}

Espira de radio $R$ en el plano $xy$, con corriente $I$; punto de observación
sobre el eje $z$ a distancia $a$.

\[
\begin{gathered}
\vec r_2 = a\,\hat k,\qquad
\vec r_1 = R\,\hat e_r = R\left(\cos\theta\,\hat\imath + \sin\theta\,\hat\jmath\right),\\
d\vec l_1 = R\,d\theta\,\hat e_\theta = R\,d\theta\left(-\sin\theta\,\hat\imath + \cos\theta\,\hat\jmath\right).
\end{gathered}
\]

\begin{figure}[H]
  \centering
  % Clase 17 — espira circular de radio R en el plano xy; punto de observación
% sobre el eje z a distancia a. Vista en perspectiva. Como r_2 es ortogonal al
% plano, |r_2 - r_1| = sqrt(a^2 + R^2) para todo punto de la espira.
\begin{tikzpicture}[scale=1]

  % ejes (perspectiva)
  \draw[figaxis] (0,0) -- (-2.3,-1.25);
  \node[figlbl] at (-2.5,-1.4) {$x$};
  \draw[figaxis] (0,0) -- (3.6,0);
  \node[figlbl] at (3.85,0) {$y$};
  \draw[figaxis] (0,0) -- (0,3.9);
  \node[figlbl] at (0,4.15) {$z$};

  % espira
  \draw[figblue, line width=1.3pt] (0,0) ellipse (2.4 and 0.8);
  \draw[figvecres] (-0.25,-0.8) -- (0.35,-0.8);
  \node[figlbl, text=figred] at (0.05,-1.12) {$I$};

  % r_1 y dl_1 (punto de la espira en el ángulo de 325 grados de la elipse)
  \coordinate (P) at (1.966,-0.459);
  \draw[figvec] (0,0) -- (P);
  \node[figlbl] at (0.95,-0.47) {$\vec r_1$};
  \draw[figvecres] (P) -- ++(0.62,0.3);
  \node[figlbl, text=figred] at (2.95,-0.25) {$d\vec l_1$};

  % punto de observacion
  \coordinate (Q) at (0,2.9);
  \draw[figaux, -{Latex[length=2mm,width=1.4mm]}] (P) -- (Q);
  \node[figlbl] at (1.75,1.5) {$\vec r_2-\vec r_1$};
  \fill[figblue] (Q) circle (1.6pt);
  \node[figlbl, anchor=east] at (-0.15,1.5) {$a$};
  \draw[figvecres] (Q) -- ++(0,0.75);
  \node[figlbl, text=figred, anchor=east] at (-0.12,3.4) {$\vec B$};

  \node[figlblaux] at (-0.75,0.5) {$R$};
  \draw[figaux] (0,0) -- (-1.7,0.565);

\end{tikzpicture}

  \caption{Espira de radio $R$ en el plano $xy$ y punto de observación sobre el eje. La distancia $|\vec r_2-\vec r_1|$ es la misma para todos los puntos de la espira, y las componentes del campo perpendiculares al eje se cancelan.}
\end{figure}

Producto vectorial, término a término:

\[
d\vec l_1\times(\vec r_2-\vec r_1)
= R\,d\theta\left(-\sin\theta\,\hat\imath+\cos\theta\,\hat\jmath\right)\times
\left(a\,\hat k - R\cos\theta\,\hat\imath - R\sin\theta\,\hat\jmath\right)
\]

\begin{center}
\renewcommand{\arraystretch}{1.35}
\small
\begin{tabularx}{\textwidth}{>{\raggedright\arraybackslash}X>{\raggedright\arraybackslash}X}
\toprule
\textbf{Producto} & \textbf{Resultado} \\
\midrule
$-\sin\theta\,\hat\imath\times a\,\hat k$ & $a\sin\theta\,\hat\jmath$ \\
$-\sin\theta\,\hat\imath\times(-R\sin\theta\,\hat\jmath)$ & $R\sin^2\theta\,\hat k$ \\
$\cos\theta\,\hat\jmath\times a\,\hat k$ & $a\cos\theta\,\hat\imath$ \\
$\cos\theta\,\hat\jmath\times(-R\cos\theta\,\hat\imath)$ & $R\cos^2\theta\,\hat k$ \\
\bottomrule
\end{tabularx}
\end{center}

\[
d\vec l_1\times(\vec r_2-\vec r_1) = R\,d\theta\left[
a\left(\cos\theta\,\hat\imath + \sin\theta\,\hat\jmath\right) + R\,\hat k\right].
\]

Como $\vec r_2$ es ortogonal al plano de la espira, por Pitágoras
$|\vec r_2-\vec r_1| = \sqrt{a^2+R^2}$, constante. Las integrales de
$\cos\theta$ y $\sin\theta$ entre $0$ y $2\pi$ se anulan; el término según
$\hat k$ es constante y da $2\pi$:

\begin{keybox}
\[
\boxed{\ \vec B = \frac{\mu_0 I}{2}\,\frac{R^2}{\left(a^2+R^2\right)^{3/2}}\,\hat k\ }
\]
\end{keybox}

\subsection{Solenoide finito}

\textbf{Ejercicio propuesto} (el docente da el resultado y no hace la integral).
Solenoide de largo $L$, base circular de radio $R$, con $N$ espiras en total, o
sea $n = N/L$ espiras por unidad de longitud. \textbf{No} se lo supone infinito. La
aproximación de solenoide es que las espiras están tan pegadas que la distancia
entre ellas es despreciable; se lo trata como una suma de espiras, usando §5.3.

Con el origen en un borde y el punto de observación sobre el eje en
$\vec r_2 = z_0\,\hat k$, el campo es paralelo a $Oz$:

\begin{keybox}
\[
\boxed{\ B_z(z_0) = \frac{\mu_0\,n\,I}{2}\left[
\frac{L-z_0}{\sqrt{(L-z_0)^2+R^2}} + \frac{z_0}{\sqrt{z_0^2+R^2}}\right]\ }
\]
\end{keybox}

Si el punto está adentro ($0<z_0<L$) y \textbf{lejos de los dos bordes}
($z_0\gg R$ y $L-z_0\gg R$), cada cociente vale 1 y se recupera el solenoide
ideal:

\[
B_z = \mu_0\,n\,I .
\]

\subsection{Corrección al solenoide ideal}

Se puede ir un poco más lejos y estimar el error. Con $z_0\gg R$, por Taylor:

\[
\frac{z_0}{\sqrt{z_0^2+R^2}} = \left(1+\frac{R^2}{z_0^2}\right)^{-1/2}
= 1 - \frac12\,\frac{R^2}{z_0^2} + \mathcal O\!\left(\frac{R^4}{z_0^4}\right),
\]
y lo mismo para el otro término con $z_0\to L-z_0$ (es el mismo cálculo). El
parámetro de desarrollo ya es cuadrático en $R$, así que el orden siguiente es
cuártico. Reemplazando:

\begin{keybox}
\[
\boxed{\ B_z(z_0) \simeq \mu_0\,n\,I\left[1 - \frac14\,\frac{R^2}{z_0^2}
- \frac14\,\frac{R^2}{(L-z_0)^2}\right]\ }
\]
\end{keybox}

El primer término es el solenoide ideal; los otros dos, las primeras
correcciones. Muestran \textbf{cuándo la aproximación se vuelve mala}: al acercarse
a un borde ($z_0\to0$) o al otro ($L-z_0\to0$) la corrección explota.

\begin{figure}[H]
  \centering
  % Clase 17 — campo sobre el eje de un solenoide finito, en unidades del
% solenoide ideal mu_0 n I. Fórmula dada en clase, con L = 10 R:
%   B_z / (mu_0 n I) = (1/2) [ (L - z0)/sqrt((L - z0)^2 + R^2) + z0/sqrt(z0^2 + R^2) ].
% Las líneas verticales marcan los bordes del solenoide (z0 = 0 y z0 = L).
% En el centro vale 0,98: el 2 % de corrección calculado en clase.
\begin{tikzpicture}
  \begin{axis}[
    emfig, width=12cm, height=6cm,
    xmin=-4, xmax=14, ymin=0, ymax=1.15,
    xtick={-4,-2,0,2,4,6,8,10,12,14}, ytick={0,0.5,1},
    xlabel={$z_0/R$}, ylabel={$B_z/(\mu_0 n I)$},
    legend style={at={(0.5,-0.42)}, anchor=north, legend columns=2, /tikz/every even column/.append style={column sep=1em}},
  ]
    \addplot[figaux, forget plot] coordinates {(0,0) (0,1.15)};
    \addplot[figaux, forget plot] coordinates {(10,0) (10,1.15)};
    \addplot[figred, dashed, domain=-4:14, samples=2] {1};
    \addlegendentry{solenoide ideal}
    \addplot[figblue, domain=-4:14, samples=181]
      {0.5*((10-x)/sqrt((10-x)^2+1) + x/sqrt(x^2+1))};
    \addlegendentry{solenoide finito, $L = 10R$}
  \end{axis}
\end{tikzpicture}

  \caption{Campo sobre el eje de un solenoide de largo $L = 10R$, según la fórmula dada en clase, en unidades del solenoide ideal. Las líneas verticales marcan los bordes. En el centro vale $0{,}98$; en cada borde, la mitad.}
\end{figure}

\begin{notebox}[Con números]
En el punto medio, $z_0 = L/2$, de un solenoide «bastante
alargado», $L = 10R$: cada corrección es del 1 \%, y el error total es del
orden del \textbf{2 \%}.
\end{notebox}

\noindent\emph{La próxima clase: primero la delta de Dirac, postergada desde el principio del
curso, que permite vincular distribuciones continuas con cargas discretas; y
con ella, la ley de Ampère y su equivalencia con Biot y Savart.}

\end{document}
